% ----------------------------------------------------------
% Seção 1 — Introdução
% Na introdução do seu trabalho, apresente o tema, o problema,
% a justificativa, os objetivos e a organização do texto.
% ----------------------------------------------------------
\chapter{Introdução}\label{cap:introducao}

Este documento é um guia e, ao mesmo tempo, o modelo que ele descreve. Cada seção primária explica uma parte do trabalho acadêmico e demonstra, no próprio código-fonte, como produzi-la com a classe abn\TeX2 \cite{abntex2classe}, que estende a classe \textit{memoir} \cite{memoir} para atender às normas da Associação Brasileira de Normas Técnicas (ABNT). Para escrever o seu trabalho, leia a explicação de cada seção e substitua o texto pelo seu.

O \LaTeX\ é um sistema de preparação de documentos construído sobre o \TeX\ \cite{knuth1984texbook}. Nele, o autor descreve a estrutura do texto (seções, citações, figuras, tabelas) e o sistema cuida da formatação \cite{lamport1994}. Em trabalhos acadêmicos, isso significa que numeração de seções, sumário, listas de ilustrações, citações e referências são gerados automaticamente e permanecem consistentes a cada alteração. Para aprofundar o estudo do \LaTeX, consulte \citeonline{mittelbach2004} e \citeonline{dongen2012}.

\section{Como usar este guia}\label{sec:como-usar}

O modelo é distribuído em duas formas com o mesmo conteúdo: o modelo de arquivo único, em que todo o texto está no \texttt{main.tex}, e o modelo modular, organizado em arquivos separados:

\begin{alineas}
	\item \texttt{main.tex}: preâmbulo, dados do trabalho (título, autor, orientador), elementos pré-textuais e a ordem das seções;
	\item \texttt{secoes/}: um arquivo por seção primária (só no modelo modular);
	\item \texttt{pos-textuais/}: apêndices e anexos (só no modelo modular);
	\item \texttt{referencias.bib}: base de referências bibliográficas;
	\item \texttt{assets/}: imagens;
	\item \texttt{misc/abntex2-UFV.sty}: pacote com a capa, a folha de rosto e a folha de aprovação da UFV.
\end{alineas}

No modelo modular, para acrescentar uma seção primária, crie um arquivo em \texttt{secoes/} e inclua-o no \texttt{main.tex} com o comando \verb|\input|; no modelo de arquivo único, escreva-a diretamente no \texttt{main.tex}. A compilação está descrita na \autoref{cap:ferramentas}.

\section{Escolha da versão da ABNT}\label{sec:versao-abnt}

As normas de citação (NBR 10520) e de referências (NBR 6023) foram revistas, e alguns cursos ainda adotam as edições anteriores. Por isso, o pacote \texttt{abntex2-UFV} aceita a opção \texttt{abnt}, definida no \texttt{main.tex}:

\begin{minted}{latex}
\usepackage[abnt=2023]{misc/abntex2-UFV}  % ou abnt=2018, abnt=2002
\end{minted}

O \autoref{qua:versoes-abnt} resume cada opção. A NBR 6023 teve nova edição em 2025 \cite{abnt6023-2025}, com mudanças pontuais nas referências de coleções e fascículos de periódicos, de artigos eletrônicos com e-localização, de eventos no todo e de jurisprudência; o estilo bibliográfico usado nas opções 2018 e 2023 segue a edição de 2018, por isso confira com a edição de 2025 as referências desses tipos. Este documento foi compilado com a opção \texttt{abnt=\imprimirversaoabnt}. Antes de escolher, confirme com o orientador ou com o regulamento do curso qual edição deve ser seguida.

\begin{quadro}[htbp]
	\caption{Versões da ABNT disponíveis no modelo}\label{qua:versoes-abnt}
	\centering
	\begin{tabular}{p{0.15\textwidth}p{0.32\textwidth}p{0.35\textwidth}}
		\toprule
		\textbf{Opção} & \textbf{Citações} & \textbf{Referências} \\
		\midrule
		\texttt{2023} & NBR 10520:2023 \cite{abnt10520}: autor em maiúsculas e minúsculas & NBR 6023:2018 \cite{abnt6023} \\
		\texttt{2018} & NBR 10520:2002 \cite{abnt10520-2002}: autor em maiúsculas & NBR 6023:2018 \\
		\texttt{2002} & NBR 10520:2002: autor em maiúsculas & NBR 6023:2002 \cite{abnt6023-2002} \\
		\bottomrule
	\end{tabular}
	\fonte{Elaborado pelo próprio autor.}
\end{quadro}

\section{Objetivos}\label{sec:objetivos}

Em um trabalho acadêmico, esta seção apresenta o objetivo geral e os objetivos específicos. Os objetivos começam com verbo no infinitivo e devem ser verificáveis ao final do trabalho. Para planejar a pesquisa e redigir os objetivos, consulte \citeonline{gil2017} e \citeonline{marconi2017}.

\subsection{Objetivo geral}\label{sec:objetivo-geral}

Orientar a escrita de trabalhos acadêmicos da UFV em \LaTeX, conforme as normas da ABNT.

\subsection{Objetivos específicos}\label{sec:objetivos-especificos}

\begin{alineas}
	\item apresentar a estrutura do trabalho acadêmico e suas regras de formatação;
	\item demonstrar os tipos de citação e a elaboração de referências;
	\item demonstrar a inserção de ilustrações, tabelas, equações, algoritmos e código-fonte;
	\item descrever um fluxo de trabalho para compilar e versionar o documento.
\end{alineas}

\section{Organização do texto}\label{sec:organizacao}

A \autoref{cap:estrutura} apresenta a estrutura do trabalho acadêmico. A \autoref{cap:citacoes} trata das citações e a \autoref{cap:referencias}, das referências. A \autoref{cap:elementos} demonstra ilustrações, tabelas e demais elementos. A \autoref{cap:ferramentas} descreve as ferramentas de trabalho, e a \autoref{cap:conclusao} apresenta as considerações finais.
